\section{Tractable Reformulation}
\label{sec-reformulation}

Throughout this section, the training data, the test residuals, and $\theta\ge0$ are fixed. We use the observed residuals to partition both the residual and input spaces. These partitions express the DW restrictions as cumulative cell probabilities and reduce the Lagrangian duals to finite linear programs. A joint-distribution construction establishes exact finite primal representations, after which LP duality and the ordering of the cells yield strong duality and closed-form bounds. The main development takes $\Xi=\mathbb R$; Remark~\ref{rem-interval-support} treats known interval support.

For each empirical probability level $q_j:=j/N^{\mathrm{te}}$, define the scalar acceptance interval
\begin{equation}
  [\alpha_j(\theta),\beta_j(\theta)]
  :=\left\{p\in[0,1]:
  N^{\mathrm{te}}\phi(q_j,p)-\psi(q_j,p)\le\theta\right\},
  \qquad j=0,\ldots,N^{\mathrm{te}}.
  \label{eq-dw-rank-acceptance}
\end{equation}
Let $\xi_{(1)}^{\mathrm{te}}<\cdots<\xi_{(K)}^{\mathrm{te}}$ be the distinct ordered test residuals, and set
\begin{equation}
  c_0:=0,\quad
  c_k:=\sum_{j=1}^{N^{\mathrm{te}}}\mathbf1_{\{\xi_j^{\mathrm{te}}\le\xi_{(k)}^{\mathrm{te}}\}},\quad
  l_k:=\alpha_{c_k}(\theta),\quad
  u_k:=\beta_{c_{k-1}}(\theta),\qquad k=1,\ldots,K.
  \label{eq-dw-counts-and-levels}
\end{equation}
The counts retain all multiplicities, so $c_K=N^{\mathrm{te}}$ even when test residuals coincide. Set $\xi_{(0)}^{\mathrm{te}}:=-\infty$ and $\xi_{(K+1)}^{\mathrm{te}}:=+\infty$. Separating the observed values from the intervening intervals gives $J:=2K+1$ residual cells:
\begin{equation}
  \Xi_{2k}:=\left\{\xi_{(k)}^{\mathrm{te}}\right\}
    \quad\left(k=1,\ldots,K\right),
  \qquad
  \Xi_{2k+1}:=\left(\xi_{(k)}^{\mathrm{te}},\xi_{(k+1)}^{\mathrm{te}}\right)
    \quad\left(k=0,\ldots,K\right).
  \label{eq-residual-cells}
\end{equation}
The singleton cells distinguish probability strictly below an observed residual from probability at or below it.

Define the surrogate-response CDF by $H(t):=\mathbb P\left\{\hat g(\bm X)\le t\right\}$, with $H(-\infty):=0$ and $H(+\infty):=1$. For $i=0,\ldots,K$, the same residual levels classify inputs by the location of the failure threshold $-\hat g(\bm x)$:
\begin{equation}
  \mathcal X_i:=\left\{\bm x\in\mathcal X:
    \xi_{(i)}^{\mathrm{te}}\le-\hat g(\bm x)<\xi_{(i+1)}^{\mathrm{te}}\right\},
  \qquad
  w_i:=\mathbb P(\mathcal X_i)
    =H\left(-\xi_{(i)}^{\mathrm{te}}\right)-H\left(-\xi_{(i+1)}^{\mathrm{te}}\right).
  \label{eq-input-cells}
\end{equation}
Thus the two spaces are partitioned as
\[
  \mathcal X=\bigsqcup_{i=0}^{K}\mathcal X_i,
  \qquad
  \Xi=\bigsqcup_{j=1}^{J}\Xi_j,
  \qquad
  H\left(-\xi_{(k)}^{\mathrm{te}}\right)=\sum_{i=k}^{K}w_i,
  \quad k=1,\ldots,K.
\]
The half-open input convention assigns $-\hat g(\bm x)=\xi_{(k)}^{\mathrm{te}}$ to $\mathcal X_k$. The tail sum therefore includes threshold equalities even when $H$ has atoms.

\begin{proposition}[CDF band and residual cell probabilities]
  \label{prop-finite-rank-constraints}
  Each acceptance set in \eqref{eq-dw-rank-acceptance} is a nonempty closed interval containing $q_j$. Define
  \[
    L_\theta(t):=\alpha_{N^{\mathrm{te}}\widehat F^{\mathrm{te}}(t)}(\theta),
    \qquad
    U_\theta(t):=\beta_{N^{\mathrm{te}}\widehat F^{\mathrm{te}}(t)}(\theta).
  \]
  The DW ambiguity set has the equivalent representations
  \[
    \mathcal Q^{\mathrm{DW}}(\theta)
    =\left\{\mathbb Q\in\mathcal M_1(\Xi):
      L_\theta(t)\le F_{\mathbb Q}(t)\le U_\theta(t),\ \forall t\in\mathbb R\right\}
  \]
  and
  \begin{equation}
    \mathcal Q^{\mathrm{DW}}(\theta)
    =\left\{\mathbb Q\in\mathcal M_1(\Xi):
      \sum_{j<2k}\mathbb Q(\Xi_j)\le u_k,\quad
      \sum_{j\le2k}\mathbb Q(\Xi_j)\ge l_k,\quad k=1,\ldots,K\right\}.
    \label{eq-finite-rank-constraints}
  \end{equation}
  Equivalently, the rank probabilities satisfy
  \[
    \mathbb Q\left\{\xi<\xi_{(k)}^{\mathrm{te}}\right\}\le u_k,
    \qquad
    \mathbb Q\left\{\xi\le\xi_{(k)}^{\mathrm{te}}\right\}\ge l_k,
    \qquad k=1,\ldots,K.
  \]
  The sequences $l_k$ and $u_k$ are nondecreasing and satisfy
  \[
    0\le l_k,u_k\le1,\qquad l_{k-1}\le u_k,\qquad l_k\le u_{k+1},
    \qquad k=1,\ldots,K,
  \]
  where $l_0:=0$ and $u_{K+1}:=1$.
\end{proposition}
\begin{proof}
  For fixed $q$, the Bernoulli divergence increases as $p$ moves away from $q$, while $\psi(q,p)$ cannot increase as the interval between $q$ and $p$ expands. The corrected discrepancy therefore decreases toward $q$ from below and increases above $q$. It is continuous for $p\in(0,1)$ and equals $-\psi_0(q)\le0$ at $p=q$. The boundary conventions in Section~\ref{sec-ambiguity-set} give a nonempty closed acceptance interval also for $q\in\{0,1\}$. Scalar inversion at each empirical CDF level yields the stated CDF band, following \citet[Supplement, Section~S.6]{L.J23}.

  For $k=1,\ldots,K-1$, the empirical CDF on $\left[\xi_{(k)}^{\mathrm{te}},\xi_{(k+1)}^{\mathrm{te}}\right)$ equals $c_k/N^{\mathrm{te}}$. Monotonicity of a candidate CDF makes its lower band constraint equivalent to $F_{\mathbb Q}\left(\xi_{(k)}^{\mathrm{te}}\right)\ge\alpha_{c_k}(\theta)$ and its upper band constraint equivalent to $F_{\mathbb Q}\left(\xi_{(k+1)}^{\mathrm{te}}-\right)\le\beta_{c_k}(\theta)$. The exterior intervals give $F_{\mathbb Q}\left(\xi_{(1)}^{\mathrm{te}}-\right)\le\beta_0(\theta)$ and $F_{\mathbb Q}\left(\xi_{(K)}^{\mathrm{te}}\right)\ge\alpha_{N^{\mathrm{te}}}(\theta)$; their other band constraints are automatic because $\alpha_0(\theta)=0$ and $\beta_{N^{\mathrm{te}}}(\theta)=1$. The residual partition gives
  \[
    F_{\mathbb Q}\left(\xi_{(k)}^{\mathrm{te}}-\right)=\sum_{j<2k}\mathbb Q(\Xi_j),
    \qquad
    F_{\mathbb Q}\left(\xi_{(k)}^{\mathrm{te}}\right)=\sum_{j\le2k}\mathbb Q(\Xi_j),
  \]
  proving \eqref{eq-finite-rank-constraints} and its half-line representation.

  With $p$ fixed, the corrected discrepancy decreases as $q$ approaches $p$ from either side. For empirical levels $q<q'$, an accepted value $p<q$ at level $q'$ remains accepted at level $q$; an accepted value $p>q'$ at level $q$ remains accepted at level $q'$. Since each acceptance interval contains its empirical level, both endpoints are nondecreasing. The inequalities $\alpha_{c_k}(\theta)\le c_k/N^{\mathrm{te}}\le\beta_{c_k}(\theta)$ give the stated compatibility conditions. Singleton intervals are allowed: if $q_j=1/2$ and $\theta=0$, the interval is $\{1/2\}$ because $\psi(1/2,p)=0$.
\end{proof}

Each scalar interval can be computed by one-dimensional bisection. Proposition~\ref{prop-finite-rank-constraints} then restricts only the probabilities of the residual cells; the distribution within each cell remains free.

For a feasible joint law, define its matrix of joint cell probabilities by $\pi_{ij}:=\Pi(\mathcal X_i\times\Xi_j)$. The fixed input marginal and the residual restrictions imply $\bm\pi=(\pi_{ij})\in\mathcal C_\theta$, where
\begin{equation}
  \mathcal C_\theta:=\left\{\bm\pi\in\mathbb R_+^{(K+1)\times J}:
  \begin{array}{ll}
    \displaystyle\sum_{j=1}^{J}\pi_{ij}=w_i,&i=0,\ldots,K,\\[1mm]
    \displaystyle\sum_{i=0}^{K}\sum_{j<2k}\pi_{ij}\le u_k,&k=1,\ldots,K,\\[1mm]
    \displaystyle\sum_{i=0}^{K}\sum_{j\le2k}\pi_{ij}\ge l_k,&k=1,\ldots,K
  \end{array}\right\}.
  \label{eq-cell-mass-polytope}
\end{equation}
On each product cell, the failure indicator satisfies
\begin{equation}
  \underline h_{ij}:=\mathbf1_{\{j\le2i\}}
  \le\mathbf1_{\{\xi\le-\hat g(\bm x)\}}
  \le\overline h_{ij}:=\mathbf1_{\{j\le2i+1\}},
  \qquad (\bm x,\xi)\in\mathcal X_i\times\Xi_j.
  \label{eq-cell-costs}
\end{equation}
Residual cells with $j\le2i$ always cause failure for inputs in $\mathcal X_i$, and those with $j\ge2i+2$ are always safe. Only the overlapping interval $\Xi_{2i+1}$ has different upper and lower coefficients. Thus the cell matrix records the rank constraints exactly, while the failure objective requires the constructions below.

The upper problem remains a linear program over joint measures, with the full input marginal fixed:
\begin{equation}
  \begin{aligned}
    \overline p_{\hat f}
    =\sup_{\Pi\in\mathcal M_+(\mathcal X\times\Xi)}\quad&
      \int_{\mathcal X\times\Xi}\mathbf1_{\{\xi\le-\hat g(\bm x)\}}\,\Pi(\mathrm d\bm x,\mathrm d\xi)\\
    \text{s.t.}\quad&
      \int_{\mathcal X\times\Xi}\mathbf1_{\mathcal U}(\bm x)\,\Pi(\mathrm d\bm x,\mathrm d\xi)
        =\mathbb P(\bm X\in\mathcal U),&&\forall\mathcal U\in\mathcal B(\mathcal X),\\
    &\int_{\mathcal X\times\Xi}\mathbf1_{\{\xi<\xi_{(k)}^{\mathrm{te}}\}}\,\Pi(\mathrm d\bm x,\mathrm d\xi)\le u_k,
        &&k=1,\ldots,K,\\
    &\int_{\mathcal X\times\Xi}\mathbf1_{\{\xi\le\xi_{(k)}^{\mathrm{te}}\}}\,\Pi(\mathrm d\bm x,\mathrm d\xi)\ge l_k,
        &&k=1,\ldots,K.
  \end{aligned}
  \label{eq-measure-upper}
\end{equation}
The lower problem minimizes the same integral under the same constraints. The marginal equalities fix both the input law and total mass one. The product of $\mathbb P$ and the empirical residual law is feasible for every $\theta\ge0$.

\begin{lemma}[Finite-dimensional Lagrangian duals and weak duality]
  \label{lem-weak-duality}
  Write $\bm u=(u_1,\ldots,u_K)^\top$ and $\bm l=(l_1,\ldots,l_K)^\top$. The nonnegative multipliers $\bm\lambda,\bm\gamma\in\mathbb R_+^K$ correspond to the upper and lower rank constraints. The input marginal equality has an unrestricted multiplier $\eta\in\mathcal F_{\mathrm b}(\mathcal X)$, the space of bounded Borel measurable real-valued functions. Define the residual penalty
  \begin{equation}
    \zeta_{\bm\lambda,\bm\gamma}(\xi)
    :=\sum_{k=1}^{K}\lambda_k\mathbf1_{\{\xi<\xi_{(k)}^{\mathrm{te}}\}}
      -\sum_{k=1}^{K}\gamma_k\mathbf1_{\{\xi\le\xi_{(k)}^{\mathrm{te}}\}},
    \label{eq-dual-step-function}
  \end{equation}
  whose value on $\Xi_j$ is
  \begin{equation}
    \zeta_j:=\sum_{k=1}^{K}\lambda_k\mathbf1_{\{j<2k\}}
      -\sum_{k=1}^{K}\gamma_k\mathbf1_{\{j\le2k\}},
    \qquad j=1,\ldots,J.
    \label{eq-dual-cell-values}
  \end{equation}
  The upper and lower Lagrangian dual values are given by the finite LPs
  \begin{equation}
    \begin{aligned}
      \overline d_\theta
      :=\inf_{\bm\lambda,\bm\gamma\ge\bm0,\ \bm z\in\mathbb R^{K+1}}\quad&
        \bm\lambda^\top\bm u-\bm\gamma^\top\bm l+\sum_{i=0}^{K}w_i z_i\\
      \text{s.t.}\quad&z_i+\zeta_j\ge\overline h_{ij},
        \qquad i=0,\ldots,K,\quad j=1,\ldots,J,
    \end{aligned}
    \label{eq-finite-upper-dual}
  \end{equation}
  and
  \begin{equation}
    \begin{aligned}
      \underline d_\theta
      :=\sup_{\bm\lambda,\bm\gamma\ge\bm0,\ \bm z\in\mathbb R^{K+1}}\quad&
        -\bm\lambda^\top\bm u+\bm\gamma^\top\bm l+\sum_{i=0}^{K}w_i z_i\\
      \text{s.t.}\quad&z_i-\zeta_j\le\underline h_{ij},
        \qquad i=0,\ldots,K,\quad j=1,\ldots,J.
    \end{aligned}
    \label{eq-finite-lower-dual}
  \end{equation}
  Their optimal values satisfy
  \begin{equation}
    \underline d_\theta\le\underline p_{\hat f}
      \le\overline p_{\hat f}\le\overline d_\theta.
    \label{eq-weak-duality}
  \end{equation}
\end{lemma}
\begin{proof}
  The input marginal equality is equivalent to
  \[
    \int_{\mathcal X\times\Xi}\eta(\bm x)\,\Pi(\mathrm d\bm x,\mathrm d\xi)
      =\mathbb E_{\mathbb P}\left[\eta(\bm X)\right],
    \qquad\forall\eta\in\mathcal F_{\mathrm b}(\mathcal X).
  \]
  Relaxing this equality and the rank inequalities gives the upper Lagrangian \citep[Section~5.1]{B.V04}
  \[
    \begin{aligned}
      \mathcal L^+(\Pi,\eta,\bm\lambda,\bm\gamma)
      ={}&\int_{\mathcal X\times\Xi}\mathbf1_{\{\xi\le-\hat g(\bm x)\}}\,\Pi(\mathrm d\bm x,\mathrm d\xi)
        +\mathbb E_{\mathbb P}\left[\eta(\bm X)\right]
        -\int_{\mathcal X\times\Xi}\eta(\bm x)\,\Pi(\mathrm d\bm x,\mathrm d\xi)\\
      &+\sum_{k=1}^{K}\lambda_k\left(u_k-\int_{\mathcal X\times\Xi}
          \mathbf1_{\{\xi<\xi_{(k)}^{\mathrm{te}}\}}\,\Pi(\mathrm d\bm x,\mathrm d\xi)\right)\\
      &+\sum_{k=1}^{K}\gamma_k\left(\int_{\mathcal X\times\Xi}
          \mathbf1_{\{\xi\le\xi_{(k)}^{\mathrm{te}}\}}\,\Pi(\mathrm d\bm x,\mathrm d\xi)-l_k\right)\\
      ={}&\mathbb E_{\mathbb P}\left[\eta(\bm X)\right]
        +\bm\lambda^\top\bm u-\bm\gamma^\top\bm l\\
      &+\int_{\mathcal X\times\Xi}\left[
          \mathbf1_{\{\xi\le-\hat g(\bm x)\}}-\eta(\bm x)-\zeta_{\bm\lambda,\bm\gamma}(\xi)
        \right]\,\Pi(\mathrm d\bm x,\mathrm d\xi).
    \end{aligned}
  \]
  Its supremum over finite nonnegative measures is finite exactly when the integrand is nonpositive at every $(\bm x,\xi)$. A positive value at any point permits an arbitrarily large point mass there; otherwise the supremum integral is zero, attained by the zero measure. The functional upper dual therefore minimizes $\mathbb E_{\mathbb P}[\eta(\bm X)]+\bm\lambda^\top\bm u-\bm\gamma^\top\bm l$ subject to $\eta(\bm x)\ge\mathbf1_{\{\xi\le-\hat g(\bm x)\}}-\zeta_{\bm\lambda,\bm\gamma}(\xi)$ at every pair. For the lower problem, reversing the rank penalties gives
  \[
    \mathcal L^-
      =\mathbb E_{\mathbb P}\left[\eta(\bm X)\right]
        -\bm\lambda^\top\bm u+\bm\gamma^\top\bm l
        +\int_{\mathcal X\times\Xi}\left[
          \mathbf1_{\{\xi\le-\hat g(\bm x)\}}-\eta(\bm x)+\zeta_{\bm\lambda,\bm\gamma}(\xi)
        \right]\,\Pi(\mathrm d\bm x,\mathrm d\xi).
  \]
  Its infimum over finite nonnegative measures is finite exactly when the integrand is nonnegative everywhere, in which case the infimum integral is zero. The lower functional dual maximizes $\mathbb E_{\mathbb P}[\eta(\bm X)]-\bm\lambda^\top\bm u+\bm\gamma^\top\bm l$ subject to $\eta(\bm x)\le\mathbf1_{\{\xi\le-\hat g(\bm x)\}}+\zeta_{\bm\lambda,\bm\gamma}(\xi)$ at every pair.

  For fixed rank multipliers, the best marginal functions are the pointwise envelopes. For $\bm x\in\mathcal X_i$, they satisfy the identities
  \begin{equation}
    \begin{aligned}
      \eta^+_{\bm\lambda,\bm\gamma}(\bm x)
        &:=\sup_{\xi\in\Xi}\left\{\mathbf1_{\{\xi\le-\hat g(\bm x)\}}
          -\zeta_{\bm\lambda,\bm\gamma}(\xi)\right\}
          =\max_{1\le j\le J}\left\{\overline h_{ij}-\zeta_j\right\},\\
      \eta^-_{\bm\lambda,\bm\gamma}(\bm x)
        &:=\inf_{\xi\in\Xi}\left\{\mathbf1_{\{\xi\le-\hat g(\bm x)\}}
          +\zeta_{\bm\lambda,\bm\gamma}(\xi)\right\}
          =\min_{1\le j\le J}\left\{\underline h_{ij}+\zeta_j\right\}.
    \end{aligned}
    \label{eq-dual-envelope-identities}
  \end{equation}
  To verify them, note that outside $\Xi_{2i+1}$ the two cell coefficients equal the actual failure indicator. In the overlapping interval, the lower coefficient zero is attained by $\xi=\left(\xi_{(i+1)}^{\mathrm{te}}-\hat g(\bm x)\right)/2$ if $i<K$, and by $\xi=-\hat g(\bm x)+1$ if $i=K$. The upper coefficient one is attained by $\xi=-\hat g(\bm x)$ whenever $-\hat g(\bm x)>\xi_{(i)}^{\mathrm{te}}$. The only remaining case is $-\hat g(\bm x)=\xi_{(i)}^{\mathrm{te}}$ with $i\ge1$. Then the overlapping interval is safe, but
  \[
    \zeta_{2i}=\zeta_{2i+1}-\gamma_i\le\zeta_{2i+1},
    \qquad
    1-\zeta_{2i}\ge1-\zeta_{2i+1}.
  \]
  Choosing the left endpoint therefore attains an upper-envelope value at least as large as the candidate from the overlapping interval. Together with \eqref{eq-cell-costs}, this proves both identities, including threshold equalities.

  The envelopes are bounded and constant on the Borel cells $\mathcal X_i$, so they belong to $\mathcal F_{\mathrm b}(\mathcal X)$. They respectively minimize and maximize the expectation over feasible marginal functions. Substitution yields
  \begin{align}
    \overline d_\theta
      &=\inf_{\bm\lambda,\bm\gamma\ge\bm0}\left\{
        \bm\lambda^\top\bm u-\bm\gamma^\top\bm l
        +\sum_{i=0}^{K}w_i\max_{1\le j\le J}
          \left\{\overline h_{ij}-\zeta_j\right\}\right\},
      \label{eq-upper-measure-dual}\\
    \underline d_\theta
      &=\sup_{\bm\lambda,\bm\gamma\ge\bm0}\left\{
        -\bm\lambda^\top\bm u+\bm\gamma^\top\bm l
        +\sum_{i=0}^{K}w_i\min_{1\le j\le J}
          \left\{\underline h_{ij}+\zeta_j\right\}\right\}.
      \label{eq-lower-measure-dual}
  \end{align}
  Replacing each maximum by an epigraph variable $z_i$, and each minimum by a hypograph variable $z_i$, gives \eqref{eq-finite-upper-dual} and \eqref{eq-finite-lower-dual}. This remains exact when $w_i=0$, since $z_i$ can still be chosen as the corresponding finite envelope.

  For a feasible joint law, the signs of the rank penalties give $\mathcal L^-\le p_{\hat f}(\Pi)\le\mathcal L^+$. Feasible lower dual objectives therefore bound the primal objectives from below, and feasible upper dual objectives bound them from above. Taking the respective optima over the common nonempty primal feasible set proves \eqref{eq-weak-duality}.
\end{proof}

Both duals have $3K+1$ scalar variables and $(K+1)(2K+1)$ pairwise inequalities, together with multiplier nonnegativity. The variables $z_i$ represent the optimal marginal function on each input cell. The envelope identities show that these finite LPs are the full Lagrangian duals; their derivation does not require strong duality.

\begin{theorem}[Exact finite reformulations and strong duality]
  \label{thm-finite-dual-reformulations}
  For $\mathcal Q=\mathcal Q^{\mathrm{DW}}(\theta)$ and every $\theta\ge0$, the failure probability bounds admit the exact representations
  \begin{equation}
    \overline p_{\hat f}
      =\max_{\bm\pi\in\mathcal C_\theta}\sum_{i=0}^{K}\sum_{j=1}^{J}\overline h_{ij}\pi_{ij},
    \qquad
    \underline p_{\hat f}
      =\min_{\bm\pi\in\mathcal C_\theta}\sum_{i=0}^{K}\sum_{j=1}^{J}\underline h_{ij}\pi_{ij}.
    \label{eq-finite-primal-values}
  \end{equation}
  Their Lagrangian dual values in Lemma~\ref{lem-weak-duality} satisfy
  \begin{equation}
    \overline p_{\hat f}=\overline d_\theta,
    \qquad
    \underline p_{\hat f}=\underline d_\theta.
    \label{eq-strong-duality}
  \end{equation}
  The finite primal and dual optima, the measure primal optima, and the functional dual optima are attained.
\end{theorem}
\begin{proof}
  Any feasible joint law induces a matrix in $\mathcal C_\theta$. Integrating \eqref{eq-cell-costs} therefore bounds the upper measure problem by the finite maximization from above and the lower measure problem by the finite minimization from below.

  For the reverse inequalities, fix $\bm\pi\in\mathcal C_\theta$. Draw $\bm x$ from $\mathbb P$. On $\mathcal X_i$ with $w_i>0$, select a provisional residual cell $j$ with probability $\pi_{ij}/w_i$; the row equality makes these probabilities sum to one. For $j\ne2i+1$, choose a fixed representative of $\Xi_j$, where the actual failure indicator equals both coefficients. In the overlapping cell, use the following assignments to define joint laws $\Pi^+$ and $\Pi^-$.

  For the upper problem, assign $\xi=-\hat g(\bm x)$. This causes failure and remains in $\Xi_{2i+1}$ unless $-\hat g(\bm x)=\xi_{(i)}^{\mathrm{te}}$. At equality, which requires $i\ge1$, the assignment transfers the provisional mass to the left endpoint. For any $\widetilde\xi\in\left(\xi_{(i)}^{\mathrm{te}},\xi_{(i+1)}^{\mathrm{te}}\right)$ and every observed threshold,
  \[
    \mathbf1_{\{\xi_{(i)}^{\mathrm{te}}<\xi_{(k)}^{\mathrm{te}}\}}
      =\mathbf1_{\{\widetilde\xi<\xi_{(k)}^{\mathrm{te}}\}},
    \qquad
    \mathbf1_{\{\xi_{(i)}^{\mathrm{te}}\le\xi_{(k)}^{\mathrm{te}}\}}
      \ge\mathbf1_{\{\widetilde\xi\le\xi_{(k)}^{\mathrm{te}}\}}.
  \]
  Thus every open upper probability is unchanged and every closed lower probability can only increase. The upper assignment preserves all rank inequalities and attains the coefficient $\overline h_{i,2i+1}=1$.

  For the lower problem, assign $\xi=\left(\xi_{(i+1)}^{\mathrm{te}}-\hat g(\bm x)\right)/2$ if $i<K$, and $\xi=-\hat g(\bm x)+1$ if $i=K$. These values remain in $\Xi_{2i+1}$ and strictly exceed the failure threshold, including at threshold equality. Hence the lower assignment retains the cell matrix and attains $\underline h_{i,2i+1}=0$.

  If $w_i=0$, the corresponding row is zero; set $\xi=\xi_{(1)}^{\mathrm{te}}$ on that input cell. All assignments are Borel measurable, and their finite mixtures define probability kernels. Combining them with $\mathbb P$ gives joint laws with the full prescribed input marginal, including randomization at input atoms. Both laws are feasible and satisfy
  \[
    p_{\hat f}(\Pi^+)=\sum_{i=0}^{K}\sum_{j=1}^{J}\overline h_{ij}\pi_{ij},
    \qquad
    p_{\hat f}(\Pi^-)=\sum_{i=0}^{K}\sum_{j=1}^{J}\underline h_{ij}\pi_{ij}.
  \]
  The lower law retains every provisional cell mass. The upper law may change the cell matrix at threshold equalities, but it preserves feasibility and the assigned upper objective. These identities establish the reverse inequalities required for \eqref{eq-finite-primal-values}.

  For every $\theta\ge0$, the empirical product law induces the feasible matrix
  \[
    \pi_{i,2k}^{\mathrm{emp}}:=w_i\frac{c_k-c_{k-1}}{N^{\mathrm{te}}},
    \qquad i=0,\ldots,K,\quad k=1,\ldots,K,
  \]
  with all odd columns zero. Its open and closed rank probabilities are $c_{k-1}/N^{\mathrm{te}}$ and $c_k/N^{\mathrm{te}}$, respectively, so feasibility includes $\theta=0$. The nonnegative entries of a matrix in $\mathcal C_\theta$ sum to one, and the constraints are closed. Thus $\mathcal C_\theta$ is nonempty and compact. The finite objectives lie in $[0,1]$ and attain their optima, which the joint-law constructions realize in the measure problems.

  The finite primal LPs in \eqref{eq-finite-primal-values} have row multipliers $z_i$ and rank multipliers $\lambda_k,\gamma_k$. Their dual constraints are precisely $z_i+\zeta_j\ge\overline h_{ij}$ and $z_i-\zeta_j\le\underline h_{ij}$. Finite LP strong duality therefore gives \eqref{eq-strong-duality} and dual attainment \citep[Section~5.2.3]{B.V04}. By \eqref{eq-dual-envelope-identities}, optimal rank multipliers also yield optimal bounded marginal functions for the original functional duals. Feasibility and finite optimal values suffice; strict rank inequalities are unnecessary, including at $\theta=0$.
\end{proof}

The finite formulations are exact, and their ordered failure coefficients give a further simplification. The closed-form bounds below follow from cuts between consecutive cells and admit certificates using at most one rank multiplier.

\begin{corollary}[Closed-form bounds and optimal certificates]
  \label{thm-closed-form-bounds}
  Under the conditions of Theorem~\ref{thm-finite-dual-reformulations},
  \begin{align}
    \overline p_{\hat f}
      &=\min\left\{1,\min_{k=1,\ldots,K}\left[u_k+H\left(-\xi_{(k)}^{\mathrm{te}}\right)\right]\right\},
      \label{eq-dw-closed-upper}\\
    \underline p_{\hat f}
      &=\max\left\{0,\max_{k=1,\ldots,K}\left[l_k+H\left(-\xi_{(k)}^{\mathrm{te}}\right)-1\right]\right\}.
      \label{eq-dw-closed-lower}
  \end{align}
  Each dual optimum admits a certificate with at most one nonzero rank multiplier, of value one.
\end{corollary}
\begin{proof}
  Fix nonnegative column masses $\bm y=(y_1,\ldots,y_J)^\top$ with $\sum_j y_j=1$, and write $Y_j:=\sum_{r=1}^{j}y_r$. Let $\mathcal T(\bm w,\bm y)$ be the transportation matrices with row sums $\bm w=(w_0,\ldots,w_K)^\top$ and column sums $\bm y$. This set contains the product matrix $\pi_{ij}=w_i y_j$. The residual constraints depend only on the column masses, through $Y_{2k-1}\le u_k$ and $Y_{2k}\ge l_k$.

  For fixed $\bm y$, the upper problem maximizes mass on pairs $j\le2i+1$; the lower problem maximizes safe mass on pairs $j>2i$ and subtracts it from one. Both are maximum-flow problems with source-to-row capacities $w_i$, column-to-sink capacities $y_j$, and capacity two on each permitted row-to-column edge. A transportation matrix restricts to a feasible flow on permitted pairs. Conversely, any flow can be completed to a transportation matrix because the remaining row and column masses are nonnegative with equal totals. Thus the maximum flow equals the maximum permitted mass.

  A cut of capacity one is available, so a minimum cut crosses no edge of capacity two. Its source-side columns can be taken to be exactly the neighbors of its source-side rows. Since the neighborhoods are nested, adding rows whose neighbors are already included cannot increase the capacity. Apart from empty and full row sets, upper-problem cuts therefore reduce to rows $i<k$ and columns $j<2k$, while safe-mass cuts reduce to rows $i\ge k$ and columns $j>2k$. Their capacities, for $k=1,\ldots,K$, are respectively
  \begin{align*}
    \sum_{i=k}^{K}w_i+Y_{2k-1}
      &=H\left(-\xi_{(k)}^{\mathrm{te}}\right)+Y_{2k-1},\\
    \sum_{i=0}^{k-1}w_i+1-Y_{2k}
      &=2-H\left(-\xi_{(k)}^{\mathrm{te}}\right)-Y_{2k}.
  \end{align*}
  Empty and full row sets give capacity one. The max-flow/min-cut theorem \citep{F.F56} consequently yields
  \begin{align*}
    \max_{\bm\pi\in\mathcal T(\bm w,\bm y)}\sum_{i,j}\overline h_{ij}\pi_{ij}
      &=\min\left\{1,\min_k\left[H\left(-\xi_{(k)}^{\mathrm{te}}\right)+Y_{2k-1}\right]\right\},\\
    \min_{\bm\pi\in\mathcal T(\bm w,\bm y)}\sum_{i,j}\underline h_{ij}\pi_{ij}
      &=\max\left\{0,\max_k\left[H\left(-\xi_{(k)}^{\mathrm{te}}\right)+Y_{2k}-1\right]\right\}.
  \end{align*}

  These values are nondecreasing in their respective open and closed prefixes. To attain all upper prefix limits for the maximization and all lower prefix limits for the minimization, choose
  \begin{align*}
    &y_1^+:=u_1,\qquad y_{2k}^+:=u_{k+1}-u_k,\qquad
      y_{2k+1}^+:=0,\quad k=1,\ldots,K,\\
    &y_{2k-1}^-:=0,\qquad y_{2k}^-:=l_k-l_{k-1},\quad k=1,\ldots,K,
      \qquad y_J^-:=1-l_K.
  \end{align*}
  The monotonicity and compatibility conditions in Proposition~\ref{prop-finite-rank-constraints} make both vectors nonnegative with total mass one, and
  \[
    Y_{2k-1}^+=u_k,\qquad Y_{2k}^+=u_{k+1}\ge l_k,
    \qquad Y_{2k}^-=l_k,\qquad Y_{2k-1}^-=l_{k-1}\le u_k.
  \]
  Both allocations are feasible and attain the required prefix limits. Substitution into the fixed-column values, followed by \eqref{eq-finite-primal-values}, proves \eqref{eq-dw-closed-upper} and \eqref{eq-dw-closed-lower}. The construction allows zero masses and coincident prefix limits, including at $\theta=0$.

  Finally, let $\bm e_k$ be the $k$th unit vector. Upper and lower dual certificates are respectively
  \begin{align*}
    &\bm\lambda=\bm e_k,\qquad\bm\gamma=\bm0,\qquad
      z_i=\mathbf1_{\{i\ge k\}},\\
    &\bm\lambda=\bm0,\qquad\bm\gamma=\bm e_k,\qquad
      z_i=\mathbf1_{\{i\ge k\}}-1.
  \end{align*}
  Their feasibility follows from the index inequalities
  \[
    \mathbf1_{\{i\ge k\}}+\mathbf1_{\{j<2k\}}\ge\overline h_{ij},
    \qquad
    \mathbf1_{\{i\ge k\}}-1+\mathbf1_{\{j\le2k\}}\le\underline h_{ij}.
  \]
  Their objectives are $u_k+\sum_{i=k}^{K}w_i$ and $l_k+\sum_{i=k}^{K}w_i-1$, the inner candidates in the closed-form expressions. Zero rank multipliers with $z_i=1$ or $z_i=0$ supply the outer candidates. Selecting the best candidate gives an optimal full dual solution. The envelope identities recover its functional marginal multiplier on each input cell; the tail sum includes $i=k$ and therefore uses the closed CDF value $H\left(-\xi_{(k)}^{\mathrm{te}}\right)$ also at threshold equality.
\end{proof}

